\scanpage{19}\mspage{18}
\uline{Th} Soit $\mathcal L$ un tronçon biordonné, soit $\deleted{\unsure{U}}X\subset\mathcal L$, tel qu'il existe $I\subset\mathrm{Pstrf}(\mathcal L)$ (ens des ps-tronçons fermés de $\mathcal L$), $I$ fini, \deleted{formé de} tel que
\[
\begin{aligned}
\text{a)}\quad \forall i,j\in I,\ i\ne j,\quad
\left\{\begin{aligned}
&T_i\cap T_j=\emptyset\\
&T_i\text{ et }T_j\text{ en pos. rel. mod.}\\
&\text{i.e. }\deleted{\unsure{\partial}}T_i^{\cdot}\cup \deleted{\unsure{\partial}}T_j^{\cdot}\ \text{\deleted{est} partie}\\
&\qquad\text{finie modérée (ou}\\
&\text{encore }\forall a\in\deleted{\unsure{\partial}}T_i^{\cdot},\ b\in\deleted{\unsure{\partial}}T_j^{\cdot},\\
&a\ne b,\ \text{on a}\quad\{a,b\}\in\underset{\substack{\Vert\\\mathrm{Drap}_2(\mathcal L)}}{\unsure{B|S}_{0}(\mathcal L)}
\end{aligned}\right.
\end{aligned}
\]
\deleted{\illegible{1}}
\[
\text{b)}\quad \deleted{U}X=\bigcup T_i
\]
\margin{\unsure{un|en} disj.}

Alors

$1^\circ$) $I$ est \unsure{d'unique|unique} [\unsure{$=$|.} les $T_i$ sont les comp. conn. de $U$]
\margin{NB $U=\emptyset$ sss $I=\emptyset$}
\[
\text{$2^\circ$)}\quad \deleted{\unsure{\partial U|\partial X}}\deleted{U}X^{\cdot}=\bigcup_{i\in I}T_i^{\cdot}
\]
\[
(\text{donc }\deleted{U}X^{\cdot}=\emptyset\quad\text{sss}\quad\deleted{\unsure{I=\emptyset\ \text{ou}}}\quad\deleted{U}X=L\text{ ou }\deleted{U}X=\emptyset
\]
(On dit que $\deleted{\unsure{U}}X$ est une \uline{partie} \uline{modérée}) \uline{fermée}.

\uline{Def} Si $\deleted{U}X,\deleted{\unsure{V}}Y$ sont deux parties \uline{fermées} modérées, on dit que $X$ et $Y$ sont en pos. relative modérée $X^{\cdot}$ et $Y^{\cdot}$ le sont, i.e
\[
\forall x\in X^{\cdot},\ y\in Y^{\cdot},\ \inserted{\text{avec }}a\ne y,\ \text{on a }\{a,y\}\in\mathrm{Drap}_2
\]
% margin note written sideways
\margin{\uline{Ex} Si $X$ et $Y$ sont \unsure{modérées|modérés}, avec $X\subset Y$, alors $X^{\cdot}\subset Y$ donc tt pt de $X^{\cdot}$ est en pos. rel. modérée par \unsure{$Y$} (\unsure{i.e.|lire} \unsure{p.r. à} $Y^{\cdot}$) \unsure{donc} $X$ et $Y$ \unsure{sont|sont en} en pos. modérée.}

% pale vertical watermark
\rotatebox{90}{\textsc{\unsure{ng Extra Strong}}}
